% !PRINTED 3
\weberpage{3}

Sind $a$ und $a_1$ relative Primzahlen, so ist der letzte Divisor
$a_\nu=1$. Multiplicirt man unter dieser Voraussetzung die vorstehenden
Gleichungen mit irgend einer Zahl $b$, so folgt, dass $b$ der grösste
gemeinschaftliche Theiler von $ab$ und $a_1b$ ist, und dass also jeder
gemeinschaftliche Theiler von $ab$ und $a_1$ oder von $a$ und $a_1b$
Theiler von $b$ sein muss. Ist also $a_1$ relativ prim zu $a$ und zu
$b$, so ist es auch relativ prim zu $ab$, und aus der speciellen
Annahme, dass $a_1$ eine Primzahl sei, ergiebt sich, dass ein Product
nur dann durch eine Primzahl theilbar sein kann, wenn wenigstens einer
der Factoren durch sie theilbar ist.

Ist also das Product $ab$ durch $a_1$ theilbar und ist $a_1$ relativ
prim zu $a$, so muss $b$ durch $a_1$ theilbar sein.

Sind $a,b$ irgend zwei Zahlen mit dem grössten gemeinschaftlichen
Theiler $d$ und ist $a=da'$, $b=db'$, so sind $a'$ und $b'$ relativ
prim zu einander. Jede Zahl $m$, die zugleich ein Vielfaches von $a$
und von $b$ ist, hat die Form $m=am'=da'm'$ und darin muss $a'm'$ ein
Vielfaches von $b'$ sein, also muss auch $m'$ ein Vielfaches von $b'$
sein, d.~h. jede Zahl, die zugleich ein Vielfaches von $a$ und von $b$
ist, ist durch $a'b'd$ theilbar. Diese Zahl $a'b'd$, die selbst ein
gemeinschaftliches Vielfaches von $a$ und $b$ ist, heisst daher das
\gesperrt{kleinste gemeinschaftliche Vielfache} von $a$ und $b$.

Wenn eine Zahl durch zwei relative Primzahlen theilbar ist, so ist sie
auch durch ihr Product theilbar, und wenn also eine Zahl $m$ durch
mehrere Zahlen theilbar ist, von denen je zwei zu einander relativ prim
sind, so ist sie auch durch das Product aller dieser Zahlen theilbar.

Hierauf gründet sich der Beweis des wichtigen Satzes, dass eine Zahl
$m$ immer und nur auf eine einzige Weise als ein Product von Primzahlen
dargestellt werden kann. Denn da der Theiler nicht grösser sein kann
als der Dividend, so kann $m$ sicher nur durch eine endliche Anzahl von
Primzahlen theilbar sein. Ist $a$ eine von diesen Primzahlen und
$a^\alpha$ die höchste Potenz von $a$, die in $m$ aufgeht, so ist auch
$\alpha$ eine bestimmte Zahl. Es seien nun ebenso
$b^\beta,c^\gamma,\ldots$ die höchsten Potenzen der übrigen in $m$
aufgehenden Primzahlen $b,c,\ldots$, durch die sich $m$ theilen lässt,
dann muss, da je zwei der Zahlen
$a^\alpha,b^\beta,c^\gamma,\ldots$ relativ prim sind, $m$ auch durch
das Product $a^\alpha b^\beta c^\gamma\cdots$ theilbar sein, und es
muss $m$ diesem Producte gleich sein, da sonst noch eine andere
%CONT>>
